% Irodov Part IV §4.2 original 138--149 (end of electric oscillations).
% Checked user source PDF pp.182--184, printed pp.173--175 and figures 4.34--4.36.
\electricoscproblem{138}{173}
电感 $L=0.70\,\mathrm H$、线圈内阻 $r=20\,\Omega$ 与外接无感电阻 $R$ 串联，
加有效电压 $V=220$V、角频率
$\omega=314\,\mathrm{s^{-1}}$。
求使整个电路产生的热功率最大时 $R$ 的数值与最大功率。
\exsolution
总有功电阻
$u=R+r$，感抗 $X_L=\omega L$，
有功功率
$P(u)=V^2u/[u^2+X_L^2]$。
对 $u$ 求导：
\[
\frac{\mathrm dP}{\mathrm du}
=V^2\frac{X_L^2-u^2}{(u^2+X_L^2)^2}.
\]
极大值条件 $u=X_L$。
故
\[
\boxed{R_*=\omega L-r
=314(0.70)-20\simeq0.20\,\mathrm{k\Omega}},
\]
\[
\boxed{P_{\max}=\frac{V^2}{2\omega L}
\simeq0.11\,\mathrm{kW}}.
\]
\excheck 最优 $R\ge0$ 需 $\omega L\ge r$；本题满足。题目问\emph{整个串联回路}总热功率，不能只算外接电阻部分。

\electricoscproblem{139}{173}
串联的电容和有内阻的线圈接交流电网，
调节电容使线圈产生的热功率增加
$n=1.7$ 倍。求功率因数 $\cos\varphi$ 改变百分之多少。
\exsolution
外加电源有效电压 $V$ 固定，线圈内阻为 $R$，
串联阻抗
$Z=\sqrt{R^2+X^2}$。
热功率
$P=I^2R=V^2R/Z^2
=(V^2/R)\cos^2\varphi$，
其中 $\cos\varphi=R/Z$。
两状态
$P_2/P_1=n$
使
$(\cos\varphi_2/\cos\varphi_1)^2=n$，
因功率因数取非负值，
\[
\boxed{\frac{\cos\varphi_2-\cos\varphi_1}
{\cos\varphi_1}\,100\%
=(\sqrt n-1)\,100\%\simeq30\%}.
\]
\excheck 该结论只需电源电压与线圈电阻不变；无需已知具体电容值。

\electricoscproblem{140}{173}
品质因数 $Q=100$ 的串联振荡回路，
接有效电压恒定的正弦电源。
在共振频率处热功率最大；
求将频率相对共振频率偏移百分之多少，
使有功功率降低为原来的 $1/n$，
其中 $n=2.0$。
\exsolution
弱阻尼近共振条件下
$X=\omega L-1/(\omega C)
\simeq2L(\omega-\omega_0)$。
最大功率
$P_{\max}=V^2/R$，
失谐时
$P/P_{\max}=R^2/(R^2+X^2)$。
令比值为 $1/n$，
$|X|=R\sqrt{n-1}$。
以 $Q\simeq\omega_0L/R$，
\[
\boxed{\left|\frac{\Delta\omega}{\omega_0}\right|
\simeq\frac{\sqrt{n-1}}{2Q}
=\frac1{200}=0.005=0.50\%}.
\]
\excheck 此为高Q、频偏很小的近共振公式，适用于两侧对称的近似半功率点。

\electricoscproblem{141}{173}
无感电阻 $R=0.16\,\mathrm{k\Omega}$
与具有有功内阻的线圈串联，
接有效电压 $V=220$V。
测得无感电阻两端有效电压
$V_1=80$V、线圈两端有效电压
$V_2=180$V。
求线圈内部发热功率。
\exsolution
串联电流有效值
$I=V_1/R$。
因纯电阻的电压 $\underline V_1$ 与电流同相，
线圈电压 $\underline V_2$ 相对电流滞后（电压超前） $\varphi$，
总相量
$\underline V=\underline V_1+\underline V_2$。
因此
$V^2=V_1^2+V_2^2+2V_1V_2\cos\varphi$，
线圈功率
$P_{\rm coil}=IV_2\cos\varphi
=(V_1/R)V_2\cos\varphi$。
消去 $\cos\varphi$：
\[
\boxed{P_{\rm coil}=
\frac{V^2-V_1^2-V_2^2}{2R}}
=\frac{220^2-80^2-180^2}{2(160)}
=\boxed{30\,\mathrm W}.
\]
\excheck 不能把三个电压的\emph{有效值}线性相加，
相量余弦项是关键。

\electricoscproblem{142}{173--174}
交流电源向并联的线圈和纯电阻
$R=25\,\Omega$ 供电，
总电流有效值为 $I=0.90$A，
线圈支路 $I_1=0.50$A，
电阻支路 $I_2=0.60$A。
求线圈内部热功率。
\exsolution
两支路电压相同，纯电阻支路电流与电压同相，
因此电源有效电压
$V=I_2R=15$V。
线圈电流与电压相位差为 $\varphi$，
总电流相量模满足
$I^2=I_1^2+I_2^2+2I_1I_2\cos\varphi$。
解出
$\cos\varphi=(I^2-I_1^2-I_2^2)/(2I_1I_2)$。
线圈发热
\[
P_1=VI_1\cos\varphi
=\frac R2\left(I^2-I_1^2-I_2^2\right).
\]
代数值：
\[
\boxed{P_1=\frac{25}{2}(0.90^2-0.50^2-0.60^2)
=2.5\,\mathrm W}.
\]
\excheck 无感支路的电流不能与线圈支路的电流模直接相加得到总电流。

\electricoscproblem{143}{174}
电容 $C=73\,\mathrm{\mu F}$ 与纯电阻
$R=100\,\Omega$ 并联，
通以角频率 $\omega=314\,\mathrm{s^{-1}}$ 的正弦交流电。
求整个回路的阻抗模。
\exsolution
用复数导纳
$\underline Y=1/R+i\omega C$，
模
$|Y|=\sqrt{1/R^2+\omega^2C^2}$，
因此
\[
\boxed{|Z|=\frac1{|Y|}
=\frac R{\sqrt{1+(\omega RC)^2}}}.
\]
代 $\omega RC=314(100)(73\times10^{-6})\simeq2.29$，
\[
\boxed{|Z|\simeq40\,\Omega}.
\]
\excheck 并联电容增加总导纳，故总阻抗模比单独电阻 $100\Omega$ 小。

\electricoscproblem{144}{174}
图4.35有三种交流\emph{并联}电路，
外加端子 A、B 的正弦电压为参考相位：
(a) 纯电阻 $R$ 与内阻为 $r$、自感为 $L$ 的线圈并联；
(b) 电容 $C$ 与内阻为 $R$、自感为 $L$ 的线圈并联；
(c) 纯电阻 $R$ 与串联 LC 支路并联。
各器件参数均满足\emph{总电流滞后}电压角度 $\varphi$。
给出各支路与总电流相量的图示构造。
\exsolution
以外加电压相量 $\underline V=V$ 为正实轴。
(a) 电阻电流
$\underline I_R=V/R$ 沿正实轴，
线圈电流
$\underline I_{Lr}=V/(r+i\omega L)$ 位于第四象限，
总电流
\[
\boxed{\underline I_0=\frac VR+\frac V{r+i\omega L}}.
\]
(b) 电容电流
$\underline I_C=i\omega CV$ 沿正虚轴，
线圈电流
$\underline I_{LR}=V/(R+i\omega L)$ 位于第四象限，
$\boxed{\underline I_0=\underline I_C+\underline I_{LR}}$；
题设总电流滞后要求线圈负虚部的模大于电容正虚部。
(c) 电阻电流在正实轴，
串联 LC 分支电流
$\underline I_{LC}=V/[i(\omega L-1/(\omega C))]$。
要使总电流滞后，必须
$\omega L>1/(\omega C)$，
从而 $\underline I_{LC}$ 沿负虚轴，
\[
\boxed{\underline I_0=
\frac VR-\frac{iV}{\omega L-1/(\omega C)}}.
\]
\excheck 三图都用矢量加法\emph{按头尾相接}得总电流；
(a)(b) 中线圈电流同时具有实、负虚分量，
(c) 理想 LC 支路电流纯无功。

\electricoscproblem{145}{174}
电容 $C=1.0\,\mathrm{\mu F}$ 与
具有有功电阻 $R=0.10\,\Omega$、
自感 $L=1.0\,\mathrm{mH}$ 的线圈并联，
接有效电压 $V=31$V 的交流源。
求：(a) 并联电流谐振角频率；
(b) 谐振时供电总电流、线圈支路电流和电容支路电流的有效值。
\exsolution
总导纳
$Y=i\omega C+1/(R+i\omega L)$。
电流谐振要求其虚部为零：
$\omega C-\omega L/(R^2+\omega^2L^2)=0$，
故
\[
\boxed{\omega_{\rm res}
=\sqrt{\frac1{LC}-\frac{R^2}{L^2}}
\simeq3.0\times10^4\,\mathrm{rad/s}}.
\]
此时
$R^2+\omega_{\rm res}^2L^2=L/C$，
总导纳实部
$G=R/(L/C)=RC/L$，
总供电电流
\[
\boxed{I_{\rm supply}=VRC/L
\simeq3.1\,\mathrm{mA}}.
\]
线圈支路
$\boxed{I_{\rm coil}
=V/\sqrt{L/C}=V\sqrt{C/L}\simeq0.98\,\mathrm A}$；
电容支路
$\boxed{I_C=\omega_{\rm res}CV\simeq0.98\,\mathrm A}$。
\excheck 两支路电流很大却近于反相，导致电源总电流仅几毫安；
书末各处将其约写为 $3$mA、$1$A。

\electricoscproblem{146}{174}
电容 $C$ 与有功电阻为 $R$、自感为 $L$ 的线圈并联，
交流源角频率为 $\omega$。
求总供电电流相对电源电压的相位差。
\exsolution
以源电压沿复平面实轴，
总导纳
\[
\underline Y=i\omega C+\frac1{R+i\omega L}
=\frac R{R^2+\omega^2L^2}
+i\left[\omega C-\frac{\omega L}{R^2+\omega^2L^2}\right].
\]
电流相位即 $\arg Y$，
因此
\[
\boxed{\tan\varphi=
\frac{\omega C(R^2+\omega^2L^2)-\omega L}{R}}.
\]
\excheck $\varphi>0$ 时总电流超前电压（容性）；
$\varphi<0$ 时滞后电压（感性）；$\varphi=0$ 为并联谐振。

\electricoscproblem{147}{174}
有功电阻为 $R$、自感为 $L$ 的线圈
与电容 $C$ 并联，
求交流角频率 $\omega$ 下整个回路的阻抗模。
\exsolution
由并联导纳
$Y=i\omega C+1/(R+i\omega L)$，
可写
\[
Y=\frac{1+i\omega C(R+i\omega L)}
{R+i\omega L}
=\frac{(1-\omega^2LC)+i\omega CR}
{R+i\omega L}.
\]
故
\[
\boxed{|Z|=\frac1{|Y|}
=\sqrt{\frac{R^2+\omega^2L^2}
{(1-\omega^2LC)^2+(\omega CR)^2}}}.
\]
\excheck $\omega\to0$ 时 $Z\to R$，
符合电容开路、线圈仅剩有功电阻的极限。

\electricoscproblem{148}{174}
细金属线环有电阻 $R$、自感 $L$，
在均匀外磁场中以恒定角速度 $\omega$ 转动，
外磁通随时间
$\Phi(t)=\Phi_0\cos\omega t$。
证明：(a) 感应电流的稳态规律；
(b) 维持转动所需外力平均机械功率。
\exsolution
法拉第电动势
$\mathcal E=-\dot\Phi=\omega\Phi_0\sin\omega t$。
回路
$L\dot I+RI=\mathcal E$。
其稳态解写为
\[
\boxed{I=I_m\sin(\omega t-\varphi)},\qquad
\boxed{I_m=\frac{\omega\Phi_0}
{\sqrt{R^2+\omega^2L^2}}},
\quad
\boxed{\tan\varphi=\frac{\omega L}{R}}.
\]
外力维持匀速转动所做的机械功在稳态全部补偿电阻损耗；
一周期平均
$\langle P_{\rm mech}\rangle
=R\langle I^2\rangle=RI_m^2/2$。
故
\[
\boxed{\langle P_{\rm mech}\rangle
=\frac{\omega^2\Phi_0^2R}{2(R^2+\omega^2L^2)}}.
\]
\excheck 线圈电感可在瞬时交换储能，但一周期首末磁场能相同，平均机械输入等于焦耳热。

\electricoscproblem{149}{175}
图4.36：木芯上沿轴线放两只线圈。
线圈1的电流
$I_1(t)=I_0\cos\omega t$；
线圈2短路，其有功电阻 $R$、自感 $L_2$，
互感 $M=M(x)$ 随两线圈间距 $x$ 改变。
求在稳态下两线圈间相互作用力沿 $x$ 的周期平均值。
\exsolution
线圈2的法拉第方程
$L_2\dot I_2+RI_2=-M(x)\dot I_1$，
以 $x$ 在一个振荡周期内不变近似。
复相量为
$\underline I_1=I_0$，
$\underline I_2=
-i\omega M I_0/(R+i\omega L_2)$，
所以
$\Re\underline I_2=
-\omega^2M L_2 I_0/(R^2+\omega^2L_2^2)$。
耦合磁通导致瞬时轴向力
$F_x=I_1I_2\,\mathrm dM/\mathrm dx$
（取正 $x$ 为两线圈拉开方向），
故周期平均
\[
\boxed{\langle F_x\rangle=
-\frac{\omega^2L_2I_0^2}{2(R^2+\omega^2L_2^2)}
\,M(x)\frac{\mathrm dM}{\mathrm dx}}.
\]
\excheck 对同绕向线圈，互感随距离增加通常减小（$M>0$、$M'<0$），所以
$\langle F_x\rangle>0$ 表示感生电流产生平均排斥力。
书后简答的符号取决于所选 $x$ 方向与互感有向约定，须一并指定。
